\subsection{NN-G029: an existence-preserving degree cutoff makes the unrestricted optimum decidable}
\paragraph{Statement.}
Fix a finite based trace algebra with bar-invariant Laurent
structure constants $c_{ij}^k(q)\in\mathbb Z[q^{\pm1}]$,
including a specified unit. An effective symmetric dimension
lift consists of bar-invariant $d_{ij}^k(q)\in\mathbb N[q^{\pm1}]$
that are associative and unital, dominate $|c_{ij}^k|$
coefficientwise, and have the same coefficients modulo two.
Extra positions with zero trace are permitted unless explicitly
forbidden. Its cancelling-pair cost is
\[
 C(d)=\frac12\sum_{i,j,k,e}
             \bigl(d_{ij,e}^k-|c_{ij,e}^k|\bigr).
\]
Suppose one such lift of cost at most $C_0$ is known.
Put
\[
 N_0=\sum_{i,j,k,e}|c_{ij,e}^k|,\qquad N=N_0+2C_0,
 \qquad K=\max\bigl(1,\max\{|e|:c_{ij,e}^k\ne0\}\bigr),
\]
Let $E$ consist of the positive trace exponents together with one,
put $h=|E|$, and set $n_* =\lfloor N/2\rfloor+h$.
Every lift of cost at most $C_0$ has another lift of exactly
its cost, with the same coefficient lists and the same nonzero
structure-constant positions, whose Laurent degrees are at most
\[
 B=n_*(K+1)4^{\,n_*-1}.                              \tag{G029.1}
\]
Consequently a finite terminating search computes the true
minimum over all symmetric effective associative lifts in
this presentation, without a closed-table or constant-ratio
restriction.

This is an existence-preserving compression statement.
It does \emph{not} bound the degrees of every existing lift,
assert a practical search algorithm, or identify the numerical
unrestricted optimum.

For the supplied rank-three trace table, $N_0=4443$ and $K=5$.
The NN-G027 incumbent has $C_0=3352$, so $N=11147$ and
$n_*=5578$. Every competing lift therefore has an equal-score
representative of Laurent degree
\[
 \le33468\,4^{5577}<2^{11170}.                       \tag{G029.2}
\]
One may search all $10^3$ structure-constant positions, including
new even coefficients at positions of zero trace, or retain any
specified forbidden positions. Both variants are finite and
complete for their stated presentation.

\paragraph{Elementary small-integer-point lemma.}
Let $P\subset\mathbb R^n$ be a nonempty polyhedron defined by
finitely many integer inequalities $a_r\cdot x\ge b_r$.
Equalities may be represented by their two inequalities.
Assume $P\subseteq\{x_i\ge1\}$, every row has Euclidean
norm at most four, and $|b_r|\le K$ with $K\ge1$.
If $P$ contains an integer point, it contains one satisfying
\[
 1\le x_i\le n(K+1)4^{n-1}\quad\text{for every }i.
 \tag{G029.3}
\]

Here is a complete polyhedral proof. Include the inequalities
$x_i\ge1$ in the presentation and homogenize:
\[
 \mathcal C=\{(x,t):t\ge0,\ a_r\cdot x\ge b_rt\}.
\]
Because $x_i\ge t\ge0$, this cone lies in the nonnegative
orthant and has no nonzero line. Its section
$\sum_i x_i+t=1$ is a compact polytope. Every point of
this section is a convex combination of its vertices:
if its smallest face has positive dimension, choose a
nonzero direction in that face's affine span, follow it
in both directions to the two boundary endpoints of
the compact face, and write the point as their convex
combination. The endpoints lie in faces of smaller
dimension, so induction terminates. There are only
finitely many vertices, since each is determined by
a full-rank subset of the finite active constraints.
Thus $\mathcal C$ is generated by finitely many extreme
rays.

An extreme ray with $t>0$, rescaled to $t=1$, has the
form $(v,1)$ with $v$ a vertex of $P$: a nontrivial
segment through $v$ in $P$ would contradict extremality
of that ray. An extreme ray with $t=0$ has the form
$(r,0)$ with $r$ an extreme ray of the recession cone
\[
 \operatorname{rec}(P)=\{r:a_r\cdot r\ge0\}.
\]
In particular $r_i\ge0$. Applying the cone decomposition
to an existing integer point $(x,1)$ gives
\[
 x=\sum_j\lambda_jv_j+\sum_s\mu_s r_s,
 \qquad \lambda_j,\mu_s\ge0,\quad\sum_j\lambda_j=1.
 \tag{G029.4}
\]
This proves the needed vertex-plus-recession decomposition
also when $P$ has lower affine dimension. It explains
explicitly the rays on the $t=0$ face; none is lost.

A vertex $v$ has $n$ linearly independent active integer
rows. Otherwise a small positive and negative displacement
in their common kernel would remain feasible, contradicting
that it is a vertex. Orthogonalizing the rows by Gram--Schmidt shows that the
absolute determinant is the product of the orthogonalized
row lengths, each at most its original length; this proves
the determinant bound used below. Cramer's rule then gives
\[
 |v_i|\le nK4^{n-1}:                                  \tag{G029.5}
\]
the denominator is a nonzero integer of absolute value at
least one, and each numerator is a sum of $n$ numbers
$b_r$ times cofactors bounded by $4^{n-1}$.

Likewise an extreme recession ray has $n-1$ independent
active homogeneous rows. Otherwise its active common
kernel would have a direction transverse to the ray,
and small displacements would split it inside the cone.
Its integer cofactor vector spans the ray. Choose the
nonnegative orientation and divide by the greatest common
divisor of its entries to obtain a primitive integer
ray $r$, with
\[
 0\le r_i\le4^{n-1}.                                 \tag{G029.6}
\]
The assertion also covers $n=1$, using the empty determinant
one. If the recession cone is zero there are simply no
such rays.

The ray part of (G029.4) needs at most $n$ rays. If more
are used, they are linearly dependent. A nonzero linear
relation among nonnegative nonzero rays has coefficients
of both signs. Subtract a suitable multiple of that
relation from the nonnegative coefficients until one
coefficient becomes zero, preserving all the others'
nonnegativity, and repeat. The vertex convex combination
is left unchanged.

Now replace each of these at most $n$ coefficients $\mu_s$
by its fractional part. Equivalently set
\[
 x'=x-\sum_s\lfloor\mu_s\rfloor r_s.
\]
This is an integer vector, because $x$ and the chosen
primitive rays are integral. It remains in $P$, because
it is the same convex combination of vertices plus a
nonnegative recession combination. The coefficients in
that combination are less than one, so (G029.5)--(G029.6)
give
\[
 x'_i\le nK4^{n-1}+n4^{n-1}=n(K+1)4^{n-1}.
\]
Since $P\subseteq\{x_i\ge1\}$, this proves (G029.3).
Notice that no rounding of the vertex convex combination
was performed. Integer flooring was applied only to
integer recession directions.\hfill$\square$

\paragraph{Encoding every relevant exponent collision and fixed anchor.}
Take any lift $d$ of cost at most $C_0$. Its total
coefficient mass is $N_0+2C(d)\le N$. Keep every
coefficient and every structure-constant position fixed.
For each occurrence of a positive exponent in a symmetric
pair $a(q^e+q^{-e})$, introduce one positive integer
variable $x_\nu$, and replace the pair's exponents by
$x_\nu,-x_\nu$. There are at most $\lfloor N/2\rfloor$
such variables: each pair uses coefficient mass at least
two. Central coefficients need no variable. Distinct
occurrences may use separate variables even if their
original numerical exponents agree.

Adjoin one anchor variable $y_r$ for each $r\in E$ and require
\[
 y_r=r\quad(r\in E).
\]
Thus the total number of variables is at most $n_*$.
There are $h$ such variables. These are affine integer equalities with constants
bounded by $K$; they will not be rescaled. Every recession
direction has zero anchor coordinates, so the flooring
step preserves them literally.

Form the finite collection $\mathcal L$ of all linear
forms $0$, all positive and negative individual variables,
and all sums of two such signed variables. Repetitions
such as $x_\nu+x_\nu$ are allowed. Each form has coefficient
$\ell^1$ norm at most two. It includes every exponent
appearing in a product of two structure constants,
as well as every single exponent and every fixed trace
weight $\pm r$ through its anchor.

Use the original integer exponent assignment to record
the complete order pattern of these forms. If two forms
have the same value, impose their equality. If one is
larger, impose that their difference be at least one.
Also impose positivity of every variable and the fixed
anchor equations. These are finitely many integer linear
constraints, and the original assignment is an integer
solution. Each non-anchor row has $\ell^1$ norm at most
four and right side zero or one; an anchor row has norm
one and right side at most $K$ in absolute value.
Consequently the small-integer-point lemma supplies
another solution with all coordinates at most (G029.1).

This construction preserves the exact order pattern,
not only a set of selected equalities. In particular it
prevents previously different sums from colliding and
prevents a non-anchor monomial from newly occupying a
fixed trace weight. The constraints also retain equalities
between distinct occurrences of the same original weight.

\paragraph{Why the new exponent assignment is another lift of the same cost.}
Use the small solution to redefine the monomial exponents,
leaving their coefficients unchanged. Positivity of the
variables keeps each symmetric pair away from exponent
zero; their negatives still give bar invariance. The
single-exponent order constraints keep distinct monomials
within each coefficient distinct. Thus each coefficient
list, each nonzero table position, and the total mass are
unchanged.

Domination and parity at the prescribed trace weights
are unchanged, because equality or inequality with every
anchor $\pm y_r$ and with zero was preserved and the
anchors still equal the original integers $r$.
Every monomial outside the fixed trace support stays
outside it and retains its even coefficient. The specified
unit has only its unchanged constant coefficient, so
unitality is preserved.

For associativity, fix three basis elements and an output
basis element. Each side of
\[
 \sum_h d_{ij}^h d_{hk}^{\ell}
       =\sum_h d_{jk}^h d_{ih}^{\ell}
\]
is a sum of monomial products. Their numerical coefficients
are products of the original unchanged coefficients.
Their exponents are forms in $\mathcal L$. The original
identity says that the coefficient totals agree in every
exponent-collision class. Exactly the same classes occur
after compression, so every coefficient total still
agrees. This verifies all associativity identities.
The trace coefficient norm is fixed and the dimension
mass is unchanged, hence the cancelling-pair cost is
exactly the same. This proves (G029.1).

\paragraph{The finite algorithm and its exact scope.}
There are only finitely many arrays of nonnegative integer
Laurent coefficients with total mass at most $N$ and
exponents in $[-B,B]$. Enumerate them, impose bar symmetry,
the specified unit, coefficient domination and parity,
and test associativity by exact finite polynomial
arithmetic. Optional prescribed zero positions are
retained throughout. The known incumbent ensures the
accepted set is nonempty. Minimize the cost over it.

Any better unbounded-degree lift would compress to an
accepted lift of the same cost in this finite set.
Therefore the computed minimum is the unrestricted
symmetric-lift minimum for the chosen basis and unit,
not merely a degree-window minimum. The search bound is
immense: neither polynomial time nor any useful running
time estimate follows from this argument.

For the numerical specialization, NN-G011 proves the
absolute nonunit trace mass $4424$, and the nineteen
unit positions add $19$, giving $N_0=4443$.
The supplied trace coefficient table has maximum
Laurent exponent five. The checked NN-G027 lift has
mass $11128+19=11147$. Hence $n_*\le5578$ and
$n_*(K+1)4^{n_*-1}\le33468\,4^{5577}<2^{11170}$.
The small certificate
\texttt{verify\_degree\_cutoff\_constants.py} checks
these input norms and constants; it does not perform
the enormous finite search or substitute for the proof.

\paragraph{Sparse anchors give polynomial-size certificates without an incumbent.}
The anchor set need only contain the positive exponents
actually present in the trace, together with one; all
other trace coefficients are zero. With
\[
 E=\{1\}\cup\{e>0:c_{ij,e}^k\ne0\text{ for some }i,j,k\},
 \qquad h=|E|,
\]
the proof uses $h$ anchor variables, not $K$ variables.
The same small-point lemma applies with right-hand-side
bound $K=\max E$, and gives
\[
 n\le\lfloor N/2\rfloor+h,\qquad
 \log_2 B=O(N+\log K).                              \tag{G029.7}
\]
Indeed $h\le\lfloor N_0/2\rfloor+1\le N+1$ whenever
the trace norm is at most the allowed dimension mass.
Preserving comparison with every exponent in $E$ is
sufficient: a monomial at a nonzero trace weight stays
there, and a monomial outside all those weights remains
outside the nonzero trace support. Adding the weight-one
anchor fixes the marked character used in the geometric
interpretation. The rank-three numerical bound is unchanged,
since its positive trace exponents are exactly $1,\ldots,5$.

This proves a uniform NP-membership statement. The input
is a finite based trace table encoded sparsely: rank and
basis indices, integer coefficients and signed exponents
are written in binary, and a \emph{total coefficient-mass
budget $N$} is written in unary. Equivalently one may
consider an input regime in which $N$ is polynomially
bounded by the input length. The decision problem asks
whether a symmetric effective associative unital lift
of total mass at most $N$ exists. Trace-algebra identities
and normalization of its sparse coefficient lists can
be checked as part of the input verification.

No known incumbent is needed for this decision result.
If the answer is yes, apply the compression proof to
that yes-witness with the prescribed budget $N$.
A unit on rank $r$ already uses at least $2r-1$ coefficient
units, so a yes-instance has $r\le(N+1)/2$. The trace
norm must also be at most $N$; otherwise the answer is
immediately no. These observations make both the rank
and the sparse anchor count polynomially bounded.

The compressed witness has at most $N$ nonzero monomials.
Their coefficients have $O(\log N)$ bits, their table
positions have $O(\log r)$ bits, and (G029.7) gives
$O(N+\log K)$ bits per exponent. Sorting the monomials
inside each table coefficient gives a canonical sparse
encoding of polynomial length. A verifier checks total
mass, bar symmetry, the unit, domination and parity by
sparse comparisons with the input table. It checks all
associativity identities by multiplying sparse polynomials,
adding their binary exponents, and collecting equal sums.
There are polynomially many such products and terms in
$r,N$; all coefficients and exponent sums have polynomial
bitlength. Thus the full verification is polynomial in
the input and witness lengths, proving membership in NP.

This is a certificate theorem for the existence of an
effective associative lift under a mass budget. It is
not NP-hardness, a polynomial-time algorithm, or PH1 for
actual plethysm or Kronecker multiplicities. With a binary
budget that may be exponentially large, the stated
polynomial certificate bound does not follow.

\paragraph{Geometric translation and the reduced optimization gate.}
The exponent variables are characters in a common finite
presentation of a diagonalizable group. More explicitly,
start with a free character group on these symbols and
one marked symbol $e_0$. Impose the homogeneous character
relations corresponding to the exponent collisions and
the anchor relations $y_r-r e_0=0$. An integral cocharacter
to this diagonalizable group, with value one on $e_0$,
assigns the integer grading weights. The strict order
conditions select an open rational polyhedral region
in that affine cocharacter slice. The original and the
compressed solutions are two integral cocharacters in
the same finite collision region, both fixing the marked
trace weights.

SGA 3, I, 4.7.2--4.7.3 identifies these character
presentations with diagonalizable-group comodules.
Pulling back their finite multiplicity objects along
the two cocharacters preserves the displayed finite
multiplication relations; NN-G010 and archive cote 2,
scan 17 then give the corresponding finite-kernel
realizations. Thus the actual conclusion is existence
of an equal-cost realization with bounded grading
weights and the same specified trace table. No equivalence
of the two categories or canonical choice of grading
is asserted.

The reassignment is not a homomorphism of the original
whole weight lattice $\mathbb Z$ to itself, and is not
a single Veronese scaling: the weight-one anchor remains
one. It preserves the finite character relations needed
for this multiplication presentation. The unrestricted
symmetric-lift optimum is now a finite effective search
problem, and unary-mass-budget existence has polynomial-size
verifiable certificates. Its value and any practical
optimization method remain open. Nothing here computes actual GCT multiplicities,
proves PH0/PH1, or transfers a nonstandard purity statement.
